\begingroup
\setlength{\tabcolsep}{4pt}
\begin{tabularx}{\textwidth}{>{\bfseries}L{0.17\textwidth}L{0.28\textwidth}X}
\toprule
Jelölés & Jelentés & Megjegyzés \\
\midrule
$p,q$ & két különböző nagy prímszám & Ezek titkos kiinduló adatok. Gyakorlati RSA-ban nagyon nagy, véletlenszerű prímeket használnak. \\
$n=pq$ & modulus & A nyilvános és a titkos kulcsban is szerepel. A biztonság alapja, hogy $n$ faktorizálása nagy számokra nehéz. \\
$\varphi(n)$ & Euler-féle függvény értéke & Két prím szorzatára $\varphi(n)=(p-1)(q-1)$. \\
$e$ & nyilvános kitevő & Úgy kell választani, hogy $1<e<\varphi(n)$ és $\operatorname{lnko}(e,\varphi(n))=1$. \\
$d$ & titkos kitevő & $d\equiv e^{-1}\pmod{\varphi(n)}$, vagyis $ed\equiv1\pmod{\varphi(n)}$. \\
$(e,n)$ & nyilvános kulcs & Ezzel bárki titkosíthat: $C\equiv M^e\pmod n$. \\
$(d,n)$ & titkos kulcs & Ezzel a címzett visszafejt: $M\equiv C^d\pmod n$. \\
\bottomrule
\end{tabularx}
\endgroup
